\chapter{INTRODUCCIÓN}

\section{CONTEXTUALIZACIÓN}

La instalación masiva de dispositivos interconectados a través del famoso Internet de las Cosas (IoT) ha logrado transformar por completo nuestra infraestructura digital a nivel global, haciendo muchísimo más fácil la automatización en áreas tan críticas como la industria 4.0, el manejo de ciudades inteligentes, la telemedicina y el monitoreo del medio ambiente, pero lamentablemente toda esta rápida expansión también ha traído consigo un enorme problema al aumentar de forma explosiva nuestro nivel de exposición ante ciberataques que resultan ser cada vez más complejos y destructivos, una situación que se vuelve todavía más complicada de manejar si tomamos en cuenta que las redes IoT son sumamente variadas y que la gran mayoría de estos pequeños equipos sufren de fuertes limitaciones físicas en cuanto a su memoria, batería y capacidad de procesamiento, lo cual provoca que sea prácticamente imposible intentar protegerlos instalándoles los típicos sistemas de seguridad tradicionales como los cortafuegos pesados o las herramientas de inspección profunda de datos (Sama et al., 2025).

A pesar de que los Sistemas de Detección de Intrusiones (IDS) basados en aprendizaje automático han avanzado muchísimo, la realidad es que al momento de identificar anomalías en entornos de red reales se topan de frente con un enorme desafío técnico conocido como el desbalance extremo de clases, una situación que ocurre porque en el día a día la inmensa cantidad de tráfico normal junto con los ataques masivos de denegación de servicio (DDoS) terminan inundando por completo los registros, provocando que los ciberataques mucho más sofisticados y silenciosos, como las inyecciones Web, los desbordamientos de búfer o los ataques de fuerza bruta, sucedan con una frecuencia tan baja que casi quedan sepultados (Ahmed et al., 2024), lo cual lamentablemente genera un fuerte sesgo algorítmico en los clasificadores convencionales, quienes terminan engañándose a sí mismos al tratar de optimizar su exactitud global a costa de ignorar casi por completo a estas muestras minoritarias que, irónicamente, son las que terminan representando las brechas de seguridad más disruptivas y peligrosas para todo el sistema.

Para tratar de mitigar esta gran vulnerabilidad, la comunidad científica ha comenzado a explorar nuevas alternativas como las arquitecturas de aprendizaje jerárquico, también conocidas como Stacking, aunque rápidamente se dieron cuenta de que simplemente mezclar varios clasificadores no es suficiente si no se le añaden mecanismos matemáticos capaces de medir la incertidumbre que se genera en esas zonas donde los datos están tan desbalanceados, por lo que integrar herramientas como la Entropía de Shannon ha resultado ser un paso clave para lograr evaluar qué tan confundidos están los modelos base, algo que, al combinarse con técnicas mucho más avanzadas para equilibrar los datos en la frontera como lo es el Borderline-SMOTE, nos permite devolverle su verdadera importancia a esos ataques menos frecuentes sin terminar ensuciando el sistema con información sintética perjudicial (Popoola et al., 2025; Nigar \& Mustafa, 2025), abriendo así la puerta para construir un gran modelo híbrido donde la poderosa sinergia entre Random Forest, XGBoost y LightGBM, trabajando todos bajo la supervisión de un meta-clasificador superior, nos ofrezca por fin una solución que no solo es capaz de adaptarse de maravilla a diferentes situaciones, sino que además cuenta con una velocidad de respuesta increíblemente rápida y ultrabaja, haciéndola ideal para proteger los complejos entornos de las redes IoT.

\section{REALIDAD PROBLEMÁTICA}

Si miramos el panorama a nivel internacional, nos damos cuenta de que la adopción masiva de tecnologías del Internet de las Cosas (IoT) tanto en la industria como en el comercio ha terminado exponiendo a las organizaciones a un número cada vez mayor de amenazas cibernéticas, una cruda realidad que golpea especialmente a las pequeñas y medianas empresas (PYMEs) de todo el mundo, quienes enfrentan un desafío gigantesco para proteger sus entornos digitales simplemente porque carecen de sistemas de ciberseguridad que se ajusten a su verdadera capacidad tecnológica (Ciampi et al., 2021), y aunque es cierto que el aprendizaje automático nos ofrece alternativas muy prometedoras para detectar estas intrusiones, el gran problema es que la mayoría de las soluciones actuales fracasan rotundamente cuando tienen que lidiar con conjuntos de datos que están altamente desbalanceados, tal como ocurre con el famoso referente internacional CICIoT2023, donde lamentablemente los ataques más sigilosos y destructivos, como aquellos basados en la web o los de fuerza bruta, apenas logran representar menos del 1\% de todo el tráfico, provocando que pasen totalmente desapercibidos y sean ignorados por los clasificadores convencionales (Nigar \& Mustafa, 2025).

Si aterrizamos esta problemática a nivel nacional, nos encontramos con que en el Perú las micro, pequeñas y medianas empresas (MYPEs y PYMEs) son el verdadero motor de la economía, ya que representan más del 98\% de todo nuestro tejido productivo y juegan un papel absolutamente fundamental en la generación de empleo para el país. Sin embargo, esta rápida carrera hacia la digitalización y la constante incorporación de dispositivos IoT en los negocios locales se está llevando a cabo bajo un escenario de altísima vulnerabilidad cibernética, una situación que se agrava enormemente debido a que estas empresas suelen operar con presupuestos muy limitados y sufren por la casi total ausencia de herramientas para detectar amenazas que estén verdaderamente adaptadas a nuestra propia realidad tecnológica (López-García et al., 2023). Toda esta preocupante realidad queda en evidencia al revisar la literatura nacional sobre el tema, donde estudios recientes aplicados sobre el dataset CICIoT2023, como los desarrollados por Raico Gallardo (2025), demuestran claramente que, si bien los clasificadores individuales logran detectar sin mucho esfuerzo el tráfico de ataques masivos, terminan fracasando de manera drástica cuando intentan identificar las intrusiones de baja frecuencia, un gran vacío que nos confirma la urgente necesidad de empezar a diseñar modelos integrados que por fin logren resolver problemas tan complejos como el rebalanceo de clases y el análisis de la incertidumbre para poder fortalecer de verdad la ciberseguridad a nivel nacional (Vargas-Hernández et al., 2022).

Si enfocamos nuestra mirada a nivel regional, específicamente en la región de Puno y en todo el sur del Perú, nos damos cuenta de que las pequeñas y medianas empresas dedicadas a la producción, el comercio y los servicios están viviendo una paulatina pero constante automatización a través del uso de sensores y dispositivos conectados. Sin embargo, todo este importante proceso de modernización se está llevando a cabo enfrentando grandes limitaciones en la infraestructura de red y operando con recursos computacionales bastante reducidos (Quispe-Mamani et al., 2023), un escenario que nos hace ver la urgencia de empezar a implementar modelos de detección de intrusos que, al ser probados rigurosamente en un entorno de laboratorio, nos permitan evaluar de manera segura qué tan factibles y eficientes resultan estas soluciones para nuestra realidad técnica. Es justamente por esta razón que la presente investigación nace con el propósito de cubrir ese gran vacío, desarrollando un modelo híbrido altamente optimizado que logra integrar la técnica Borderline-SMOTE, la Entropía de Shannon y la arquitectura de Stacking de clasificadores dentro de un entorno de laboratorio, ofreciendo así una solución concreta que no solo nos permite elevar drásticamente la capacidad para detectar esos escurridizos ataques minoritarios, sino que además lo hace exigiendo un mínimo consumo de recursos para adaptarse a las verdaderas posibilidades tecnológicas de nuestra región.

\section{FORMULACIÓN DEL PROBLEMA}

\subsection{Problema General}
¿En qué medida el Modelo híbrido basado en Stacking, XGBoost y LightGBM mejora la eficacia en la detección de ataques minoritarios en redes IoT?

\subsection{Problemas Específicos}
\begin{itemize}
\item \textbf{PE1:} ¿En qué grado la técnica de rebalanceo Borderline-SMOTE incrementa la proporción de balance de la clase minoritaria en el dataset CICIoT2023 respecto al balance original?
\item \textbf{PE2:} ¿Cuál es la capacidad de detección del Modelo híbrido basado en Stacking, XGBoost y LightGBM en la identificación de ataques minoritarios evaluada mediante las métricas de evaluación Recall, F1-Score y Accuracy?
\item \textbf{PE3:} ¿Cuál es la eficiencia computacional del Modelo híbrido basado en Stacking, XGBoost y LightGBM propuesto en términos de tiempos de entrenamiento y latencia de inferencia frente a los modelos base Random Forest y XGBoost?
\end{itemize}

\section{JUSTIFICACIÓN}

\subsection{Justificación Técnica y Científica}
Desde el punto de vista técnico, esta investigación nace de la gran necesidad de superar las limitaciones que tienen los clasificadores individuales cuando se enfrentan a datos sumamente desbalanceados en las redes IoT, y aunque la literatura científica más reciente ya nos ha demostrado que utilizar enfoques de ensamble por Stacking resulta ser mucho mejor que usar modelos independientes para lograr una buena capacidad de generalización (Sama et al., 2025; Ahmed et al., 2024), nuestra propuesta busca dar un paso más allá para evolucionar el estado del arte al integrar la Entropía de Shannon como una medida cuantitativa real de la incertidumbre, colocándola directamente dentro del vector de meta-características (el cual está compuesto por 27 atributos que se dividen en 24 probabilidades de clase y 3 indicadores entrópicos), un detalle que resulta absolutamente clave porque le otorga al meta-modelo LightGBM la capacidad de tomar decisiones con muchísima más precisión justo en aquellas zonas donde suele generarse una mayor confusión, sumando a todo esto la aplicación estratégica de la técnica Borderline-SMOTE exclusivamente sobre el conjunto de entrenamiento para lograr fortalecer enormemente las fronteras de decisión de esas clases minoritarias tan vulnerables (como los ataques basados en Web y de fuerza bruta) sin caer en el temido error de la fuga de datos o data leakage, garantizando así que todo nuestro protocolo metodológico se mantenga completamente robusto y libre de sesgos.

\subsection{Justificación Económica y Operativa}
Si lo analizamos desde una perspectiva mucho más operativa y de laboratorio, resulta evidente que los entornos de prueba para redes IoT necesitan con urgencia validar soluciones que ofrezcan una altísima velocidad de respuesta consumiendo la menor cantidad de memoria posible, y aunque los famosos modelos basados en aprendizaje profundo (Deep Learning) son muy potentes, la cruda realidad es que exigen demasiados recursos físicos y tardan muchísimo en ejecutarse, por lo que el modelo propuesto en esta investigación logra marcar una verdadera diferencia al demostrar una eficiencia sobresaliente en estas condiciones de prueba, logrando alcanzar una latencia de inferencia de apenas $8.65\,\mu\text{s}$ por muestra (lo que equivale a procesar más de 115,000 paquetes por segundo) junto con un ligerísimo consumo de memoria RAM de tan solo $257.89\text{ MB}$, cifras contundentes que respaldan y fundamentan por completo su viabilidad técnica para poder ser implementado en el futuro sin ningún problema dentro de aquellos dispositivos que cuentan con recursos sumamente limitados.

\subsection{Justificación Social y Práctica}
En cuanto a su impacto y relevancia social, el verdadero valor de este trabajo radica en brindar evidencia tanto empírica como tecnológica que nos permita proteger a fondo la integridad y la privacidad de las redes de información que utilizan las MYPEs y PYMEs en nuestra región y en todo el país, sobre todo si tenemos en cuenta que cuando un ciberataque logra pasar desapercibido en los sistemas de estas pequeñas empresas, puede terminar provocando pérdidas económicas sumamente graves e incluso obligarlas a paralizar por completo todas sus actividades comerciales, por lo que haber logrado validar experimentalmente este modelo utilizando los 63 archivos oficiales del dataset \textbf{CICIoT2023} termina entregándole a toda la comunidad académica y técnica una herramienta de ciberseguridad verdaderamente proactiva, muy fácil de reproducir y, sobre todo, que exige un costo computacional sumamente bajo para operar de manera eficiente.

\section{OBJETIVOS}

\subsection{OBJETIVO GENERAL}
Determinar la eficacia del Modelo Híbrido basado en Stacking, XGBoost y LightGBM en la mejora de la detección de ataques minoritarios en redes IoT.

\subsection{OBJETIVOS ESPECÍFICOS}
\begin{itemize}
\item \textbf{OE1:} Evaluar la mejora en la proporción de balance de la clase minoritaria mediante la aplicación de la técnica de rebalanceo Borderline-SMOTE en el dataset CICIoT2023.
\item \textbf{OE2:} Determinar la capacidad de detección del Modelo Híbrido basado en Stacking, XGBoost y LightGBM en la identificación de ataques minoritarios utilizando las métricas de evaluación Recall, F1-Score y Accuracy.
\item \textbf{OE3:} Comparar la eficiencia computacional del modelo híbrido propuesto midiendo los tiempos de entrenamiento y latencia de inferencia frente a los modelos base Random Forest y XGBoost.
\end{itemize}

\section{HIPÓTESIS DE LA INVESTIGACIÓN}

\subsection{Hipótesis general}
El Modelo Híbrido basado en Stacking, XGBoost y LightGBM mejora significativamente la eficacia (Recall y F1-Score) en la detección de ataques minoritarios en redes IoT, superando el rendimiento de los modelos base en el dataset CICIoT2023.

\subsection{Hipótesis específicas}
\begin{itemize}
\item \textbf{HE1:} La técnica de rebalanceo Borderline-SMOTE incrementa significativamente la proporción de balance de la clase minoritaria en el dataset CICIoT2023, reduciendo el índice de desbalance y mejorando la distribución de datos en la frontera de decisión.
\item \textbf{HE2:} El Modelo Híbrido XGB-LGBM-Stacking con análisis de incertidumbre presenta una alta capacidad de detección de ataques minoritarios, evidenciando indicadores superiores en las métricas de Recall, F1-Score y Accuracy al ser validados estadísticamente mediante la prueba de rangos con signo de Wilcoxon.
\item \textbf{HE3:} El modelo híbrido propuesto optimiza la eficiencia computacional, logrando tiempos de entrenamiento y una latencia de inferencia competitivos en comparación con los modelos base Random Forest y XGBoost, garantizando su viabilidad operativa.
\end{itemize}

\section{VARIABLES}

\subsection{Variable Independiente (VI)}
Modelo híbrido basado en Stacking, XGBoost y LightGBM, optimizado con la técnica de rebalanceo de datos Borderline-SMOTE.

\subsection{Variable Dependiente (VD)}
Detección de ataques minoritarios en redes IoT, evaluada a través de la eficacia predictiva y el desempeño computacional.

\section{OPERACIONALIZACIÓN DE VARIABLES}

\begin{table}[H]
\centering
\caption{Matriz de Operacionalización de Variables de la Investigación}
\label{tab:operacionalizacion}
\small

\begin{tabularx}{\textwidth}{>{\raggedright\arraybackslash}p{2.2cm} >{\raggedright\arraybackslash}p{2.6cm} >{\raggedright\arraybackslash}p{2.6cm} >{\raggedright\arraybackslash}p{2.2cm} >{\raggedright\arraybackslash}p{2.8cm} >{\raggedright\arraybackslash}p{2.0cm}}
\toprule
\textbf{Variable} & \textbf{Definición Conceptual} & \textbf{Definición Operacional} & \textbf{Dimensiones} & \textbf{Indicadores} & \textbf{Escala} \\
\midrule
\textbf{Variable Independiente:} \newline Modelo Híbrido Stacking (RF, XGBoost, LightGBM) con Entropía de Shannon y Borderline-SMOTE & Arquitectura jerárquica de dos niveles que combina predicciones de clasificadores base (Nivel 0) con cuantificación de incertidumbre entrópica y sobremuestreo sintético en la frontera de decisión. & Pipeline computacional en Python (Conda \texttt{tesis\_iot}) que aplica Borderline-SMOTE en entrenamiento, selecciona 29 atributos, extrae 27 meta-características (24 probs + 3 entropías) y entrena un meta-LightGBM. & 1. Preprocesamiento y Rebalanceo. \newline 2. Arquitectura de Ensamble (Nivel 0 y 1). & 1. Ratio de balanceo sintético. \newline 2. Algoritmos base (RF, XGBoost GPU, LightGBM). \newline 3. Entropía de Shannon ($H$). \newline 4. Parámetros del Meta-LightGBM. & Nominal / Razón \\
\midrule
\textbf{Variable Dependiente:} \newline Detección de ataques minoritarios en redes IoT & Capacidad para identificar e inferir correctamente flujos de red anómalos de baja frecuencia (\textit{Web-based}, \textit{Brute Force}) evaluando la exactitud predictiva y la viabilidad computacional. & Evaluación experimental sobre 377,383 muestras de prueba aisladas del dataset CICIoT2023, registrando métricas de clasificación y consumo de recursos de hardware. & 1. Eficacia Predictiva. \newline 2. Desempeño Computacional. \newline 3. Significancia Estadística. & 1. Accuracy Global (\%). \newline 2. Sensitivity / Recall. \newline 3. F1-Macro / F1-Score por clase. \newline 4. Latencia de inferencia ($\mu$s/muestra). \newline 5. Uso de RAM (MB). \newline 6. $p$-valor (Prueba de Wilcoxon). & Razón / Intervalo \\
\bottomrule
\end{tabularx}
\\[5pt]
\begin{minipage}{\textwidth}
\small \textit{Nota.} Diseñado conforme a los requerimientos del modelo experimental y la rúbrica de investigación de la UNAP.
\end{minipage}
\end{table}